\documentclass[main]{subfiles}

\title[AutoML: Introduction]{Chapter 3 -- Quiz}
\subtitle{Multiple Choice Questions}
\author[Lars Kotthoff]{Lars Kotthoff}
\date{}

\titlebackground*{images/AutoML_HPO_2_2}

\newcommand{\correct}[1]{\textcolor{ForestGreen}{\textbf{#1} \checkmark}}
\newcommand{\wrong}[1]{#1}

\begin{document}

    \maketitle

%----------------------------------------------------------------------

%----------------------------------------------------------------------
\begin{frame}{Q1 \hfill \small\normalfont\textit{(single correct answer)}}

\textbf{What does HPO stand for?}

\begin{itemize}
    \item \wrong{(A) \; HyperParameter Obfuscation}
    \item \wrong{(B) \; Handy Parameter Optimization}
    \item \correct{(C) \; HyperParameter Optimization}
\end{itemize}

\end{frame}

\begin{frame}{Q2 \hfill \small\normalfont\textit{(multiple correct answers)}}

\textbf{What hyperparameters might you want to optimize?}

\begin{itemize}
    \item \correct{(A) \; Number of trees in a random forest.}
    \item \wrong{(B) \; Random seed.}
    \item \correct{(C) \; Number of layers in a neural network.}
\end{itemize}

\end{frame}

\begin{frame}{Q3 \hfill \small\normalfont\textit{(single correct answer)}}

\textbf{HPO is\ldots{}}

\begin{itemize}
    \item \wrong{(A) \; unnecessary in practice.}
    \item \correct{(B) \; often beneficial and can yield substantial performance improvements.}
    \item \wrong{(C) \; most of the time not worth the effort.}
\end{itemize}

\end{frame}

\begin{frame}{Q4 \hfill \small\normalfont\textit{(single correct answer)}}

\textbf{What is a good way to set up a configuration space?}

\begin{itemize}
    \item \wrong{(A) \; Use only default values.}
    \item \wrong{(B) \; Use largest possible range for every possible hyperparameter.}
    \item \correct{(C) \; Carefully select suitable hyperparameters and sensible ranges.}
\end{itemize}

\end{frame}

\begin{frame}{Q5 \hfill \small\normalfont\textit{(single correct answer)}}

\textbf{How can different learners be handled in a configuration space?}

\begin{itemize}
    \item \correct{(A) \; They are just another categorical hyperparameter with dependencies.}
    \item \wrong{(B) \; You have to use multiple disjoint configuration spaces.}
    \item \wrong{(C) \; They cannot be handled at all.}
\end{itemize}

\end{frame}

\begin{frame}{Q6 \hfill \small\normalfont\textit{(multiple correct answers)}}

\textbf{What are types of hyperparameters?}

\begin{itemize}
    \item \correct{(A) \; Integer.}
    \item \correct{(B) \; Categorical.}
    \item \correct{(C) \; Continuous numerical.}
\end{itemize}

\end{frame}

\begin{frame}{Q7 \hfill \small\normalfont\textit{(single correct answer)}}

\textbf{Grid search will always find the best hyperparameter configuration.}

\begin{itemize}
    \item \wrong{(A) \; Yes.}
    \item \correct{(B) \; Depends on the discretization of the grid for continuous hyperparameters.}
    \item \wrong{(C) \; No.}
\end{itemize}

\end{frame}

\begin{frame}{Q8 \hfill \small\normalfont\textit{(single correct answer)}}

\textbf{Random search often yields good results in practice.}

\begin{itemize}
    \item \correct{(A) \; Yes.}
    \item \wrong{(B) \; Only for HPO problems with many hyperparameters.}
    \item \wrong{(C) \; No.}
\end{itemize}

\end{frame}

\begin{frame}{Q9 \hfill \small\normalfont\textit{(single correct answer)}}

\textbf{Both grid search and random search can search irrelevant areas of the configuration space.}

\begin{itemize}
    \item \correct{(A) \; Yes.}
    \item \wrong{(B) \; Grid search does, but random search doesn't.}
    \item \wrong{(C) \; No.}
\end{itemize}

\end{frame}

\begin{frame}{Q10 \hfill \small\normalfont\textit{(single correct answer)}}

\textbf{Local search\ldots{}}

\begin{itemize}
    \item \wrong{(A) \; always finds the optimal hyperparameter configuration.}
    \item \correct{(B) \; can get stuck in local optima, but restarts help.}
    \item \wrong{(C) \; is less efficient than grid search.}
\end{itemize}

\end{frame}

\begin{frame}{Q11 \hfill \small\normalfont\textit{(single correct answer)}}

\textbf{What is an ``adaptive'' search?}

\begin{itemize}
    \item \correct{(A) \; A search that considers the result of previous evaluations when proposing evaluations.}
    \item \wrong{(B) \; A search that adapts itself to complex configuration spaces.}
    \item \wrong{(C) \; A search that can be adapted to multiple problems.}
\end{itemize}

\end{frame}

\begin{frame}{Q12 \hfill \small\normalfont\textit{(multiple correct answers)}}

\textbf{What are the main limitations of ablation analysis?}

\begin{itemize}
    \item \wrong{(A) \; It is not compatible with random search.}
    \item \correct{(B) \; It can only compare two hyperparameter configurations.}
    \item \correct{(C) \; It does not consider interactions between hyperparameters.}
\end{itemize}

\end{frame}

\begin{frame}{Q13 \hfill \small\normalfont\textit{(multiple correct answers)}}

\textbf{How do evolutionary algorithms change hyperparameter configurations?}

\begin{itemize}
    \item \correct{(A) \; Randomly.}
    \item \correct{(B) \; By combining two hyperparameter configurations with good performance.}
    \item \wrong{(C) \; By combining three or more hyperparameter configurations with good performance.}
\end{itemize}

\end{frame}

\begin{frame}{Q14 \hfill \small\normalfont\textit{(single correct answer)}}

\textbf{Evolutionary algorithms escape local optima by\ldots{}}

\begin{itemize}
    \item \correct{(A) \; mutating hyperparameter values.}
    \item \wrong{(B) \; random restarts.}
    \item \wrong{(C) \; They cannot escape local optima.}
\end{itemize}

\end{frame}

\begin{frame}{Q15 \hfill \small\normalfont\textit{(single correct answer)}}

\textbf{What is the main limitation of evolutionary algorithms?}

\begin{itemize}
    \item \wrong{(A) \; They cannot handle complex dependencies between hyperparameters.}
    \item \correct{(B) \; They are not straightforward to use in practice.}
    \item \wrong{(C) \; They cannot escape local optima.}
\end{itemize}

\end{frame}

\end{document}